\section{Version 1 of this study, and why it failed}
\label{app:v1}

We pre-registered a first version (commit \texttt{6a44689}) and report its failure because the
failure modes are instructive and because hiding them would misrepresent how the final design was
reached.

\paragraph{Design.} Version 1 differed in four ways:
\begin{itemize}
\item it had 25\,fps and 8 context frames;
\item restitution and friction were hidden;
\item the training pool had 16k episodes;
\item models were trained end to end: a CNN encoder and GRU with open-loop autoregressive
prediction of 24 frames, using pixel loss (REC), EMA-target latent loss (JEPA), or ground-truth
positions (STATE); 8k steps.
\end{itemize}

\paragraph{What went wrong.}
\begin{enumerate}
\item \textbf{No competence.} All three families blurred toward the mean within about ten frames.
Linear probes recovered gravity from the belief with $R^2\approx0$. Even teacher-forced
one-step predictors on ground-truth positions did not beat a constant-velocity baseline in free
flight. Gravity's per-frame signature ($g\,\Delta t^2\approx0.6$\,px) was below the models'
precision, and the one-step loss barely rewards inferring it. Knowing the exact $g$ improves the
one-step error by about $3\times10^{-6}$, against a total loss of about $2\times10^{-3}$.
\item \textbf{Memorization.} With 16k episodes, a chunk predictor's training loss kept falling
while validation error stalled. With 200k episodes, the gap closed and probe $R^2$ for $g$ rose
from 0.2 to 0.8. In this contact-rich world, generalizing physics needs far more distinct episodes
than memorizing them does.
\item \textbf{Hidden materials dominated the loss.} With the simulator we can compute how much
prediction error each unknown causes. Not knowing $(e,\mu)$ costs 7--9\,px over 24 frames; not
knowing $g$ costs 3--5\,px. Models learned contacts long before gravity.
\item \textbf{A biased gravity estimator.} Our first $g_{\rm eff}$ estimator selected free-flight
windows from the model's own predicted track. MSE-trained models hedge on the timing of bounces.
They decelerate smoothly before an uncertain contact (a ``soft bounce''), and those windows
entered the estimate as weak or even negative gravity. Restricting the fits to frames where the
\emph{reference} trajectory is in exact free flight removes the bias.
\item \textbf{A metric bug} (a missing division by the number of frames) inflated position errors
32-fold. It was caught by inspection and pinned down by a unit test.
\end{enumerate}

Diagnosing these issues required running the full version-1 evaluation on seed-0 models, which the
version-1 pilot clause forbade. Version 2 therefore uses new splits with new seeds, and nothing from
version 1 is reported as a result. Every design change was made from train/validation diagnostics
and committed in \texttt{PREREGISTRATION\_v2.md} before any version-2 evaluation.
